\appendix

\section{Anchored vector bundles and Lie algebroids}\label{preliminaries0}
In this appendix, we recall some properties of real anchored vector bundles and real Lie algebroids.
\begin{Definition}
   A {\em real anchored vector bundle $($RAVB$)$} is a pair $(A, \rho)$, where $A$ is a vector bundle over a manifold $M$ and $\rho\colon A\to TM$ is a $\R$-bundle morphism called the {\em anchor map}. We~say that a RAVB is {\em involutive} whenever $\rho(\Gamma(A))$ is a Lie subalgebra of $\Gamma(TM)$.
\end{Definition}
Given two RAVBs $(A,\rho_A)$ and $(B,\rho_B)$ over $M$ and $N$, respectively. A {\em morphism of RAVBs} is given by a bundle map $\Phi\colon A\to B$ and a map $\varphi\colon M\to N$ such that the following diagram commute:
\[\xymatrix{A\ar[d]^{\rho_A}\ar[r]^{\Phi} & B\ar[d]^{\rho_B}\\ TM\ar[r]^{T\varphi} & TN.
}\]
We recall that the Lie algebra $\mathfrak{aut}(E)$ of infinitesimal automorphism of a real vector bundle~$E$ are vector fields \smash{$\widetilde{X}\in \X(E)$} whose flow is given by automorphism of $E$. Equivalently, they are given by derivations, that is, pairs $(L,X)$, where $L\colon \Gamma(E)\to \Gamma(E)$ is a $\R$-linear operator and~${X\in \X(M)}$ satisfying $L(f\sigma)=fL(\sigma)+X(f)\sigma$. The relationship is given by the following: given~${\widetilde{X}\in \aut(E)}$, consider the pair $(L,X)$ defined by
\[
X=\pr_* \widetilde{X},\qquad \pr^*\langle \eta, X \rangle=\bigl\langle \pr^*\eta,\widetilde{X} \bigr\rangle\qquad \text{and}\qquad
L(\sigma)^{\uparrow}=\bigl[\sigma^{\uparrow},\widetilde{X}\bigr],
\]
where $\sigma\in \Gamma(E)$, $\eta\in \Omega^{1}(M)$ and $\sigma^{\uparrow}$ is the vertical lift of $\sigma$:
\begin{align*}
   \sigma\in \Gamma(E)\to \sigma^{\uparrow}\in\X(E),\qquad
 \sigma^{\uparrow}(e_m)=\frac{{\rm d}}{{\rm d}t}\bigg\rvert_{t=0}(e_m+t\sigma(m))\qquad\forall\ e_m\in E_m.
 \end{align*}
 Note that this map is injective, since $\sigma^{\uparrow}=0$, implies that $\sigma=0$ just by evaluating in the zero section of $E$.

 The {\em space of automorphisms of a RAVB $(A,\rho)$} is given by
 \[\Aut(A,\rho)=\{(\Phi, \varphi)\in \Aut(A)\mid \rho\circ\Phi=T\varphi\circ \rho\}\]
 and the {\em space of infinitesimal automorphisms of a RAVB $(A,\rho)$} is given by
 \[\mathfrak{aut}(A,\rho)=\big\{ \widetilde{X}\in \aut(A)\mid \widetilde{X}\sim_{\rho} X_T \big\},\]
 where $X_T\in \X(TM)$ is the vector field constructed with the differential of the flow of $X$. Equivalently, the infinitesimal automorphisms of $(A,\rho)$ are given by derivations $(L,X)$ satisfying~that
 \[\rho(L(\sigma))=[X, \rho(\sigma)].\]


The pullback of a RAVB $(A, \rho) $ over $M$ along the map $\varphi\colon N\to M$ is defined as the fibered product:
\[
\xymatrix{\varphi^! A\ar[d]\ar[r] & A \ar[d]^{\rho}\\
TN\ar[r]^{T\varphi} & TM.}
\]
The local description of an involutive RAVB is the following.
\begin{Theorem}[{\cite[Corollary 3.14]{bursztyn2019splitting}}]\label{splitting_ravb}
Let $(A,\rho)$ be an involutive RAVB, $m\in M$, and \smash{$N\xhookrightarrow{\iota} M$} a~submanifold containing $m$ satisfying that $\rho(A|_m)\oplus T_m N=T_m M$. Let $P=\rho(A|_m)$. Then, there exists a neighborhood $U$ of $m$ such that $A|_{U}$ is isomorphic to $\iota^! A\times TP$.
\end{Theorem}
\begin{Definition}A {\em real Lie algebroid $($RLA$)$} over a manifold $M$ is a triple $(A, [\cdot, \cdot], \rho)$, where $(A,\rho)$ is a RAVB over $M$ together with a Lie bracket on $\Gamma(A)$:
\[[\cdot,\cdot]\colon\ \Gamma(A)\times \Gamma(A)\to\Gamma(A)\]
satisfying the Leibniz property
\[[\alpha,f\beta]=f[\alpha,\beta]+\rho(\alpha)(f)\beta\]
 for all $\alpha,\beta\in \Gamma(A)$ and $f\in C^{\infty}(M)$. If the triple $(A,[\cdot,\cdot],\rho)$ satisfies all the previous conditions except the Jacobi identity, it is called a {\em skew-algebroid}, see \cite[Definition 4.1]{grabowski2022nonassociative}; they were originally introduced in \cite{kosmann1990poisson} as differential pre-Lie algebras.
\end{Definition}

Skew-algebroids $A$ are equivalent to fiber-wise linear bivectors on $A^*$, in the following way:
\begin{align*}
\{f\circ p_{A^*},g\circ p_{A^*}\}=0,\qquad
\{l_{\alpha}, f\circ p_{A^*}\}=\rho(\alpha)f,\qquad
\{l_{\alpha},l_{\beta}\}=l_{[\alpha,\beta]},
\end{align*}
where $p_{A^*}\colon A^* \to M$ is the usual bundle projection, $f, g\in C^{\infty}(M,\R)$, $\alpha,\beta\in \Gamma(A)$ and ${l_{\alpha}, l_{\beta}\in C^{\infty}(A^*,\R)}$ are defined as $l_{\alpha}(\tau)=\langle\tau,\alpha\ra$. In case these bivectors are Poisson bivectors, we obtain Lie algebroid structures.

The pullback of RLAs is defined in the same way as the pullback of RAVBs. We recall the local description of RLAs.

\begin{Theorem}[{\cite{bursztyn2019splitting, fernandes2002lie, Dufour2001NormalFF}}]
 Let $(A,[\cdot,\cdot],\rho)$ be a RLA and $m\in M$. Consider a submanifold \smash{$N\xhookrightarrow{\iota}M$} such that $\rho(A|_{m})\oplus T_m N=T_m M$ and denote by $P=\rho(A|_{m})$. Then, there exists a~neighborhood $U$ of $m$ such that $A|_{U}$ is isomorphic to $\iota^! A\times TP$.
\end{Theorem}

The following table summarizes all the algebroids used in this article:
\begin{center}
    \begin{tabular}{@{}c|c@{}}
      Real & Complex\\
      \hline
Real anchored vector bundle (RAVB): & Complex anchored vector bundle (CAVB):\\
$\R$-vector bundle and $\R$-anchor map & $\C$-vector bundle and $\C$-anchor map\\
\hline
Involutive RAVB: & Involutive CAVB:\\
A RAVB where the image of the anchor map & A CAVB where the image of the anchor map\\
 is involutive in $\Gamma(TM)$ & is involutive in $\Gamma(T_{\C}M)$ \\
\hline
 Real skew-algebroid: & Complex skew-algebroid: \\
$\R$-vector bundle, $\R$-anchor map and &
$\C$-vector bundle, $\C$-anchor map and\\
$\R$-bilinear bracket satisfying only Leibniz & $\C$-bilinear bracket satisfying only Leibniz \\
\hline
Real almost Lie algebroid: & Complex almost Lie algebroid \\
It is a real skew-algebroid where &
It is a complex skew-algebroid where\\
 the anchor map preserves the bracket & the anchor map preserves the bracket \\
\hline
Real Lie algebroid (RLA): & Complex Lie algebroid (CLA):\\
It is a real skew-algebroid where &
It is a complex skew-algebroid where\\
 the bracket satisfies Jacobi & the bracket satisfies Jacobi \\
\hline
    \end{tabular}
    \end{center}

\section{Euler-like vector fields and normal forms}\label{preliminarieslocstr}
 In this appendix we recall the properties of normal vector bundles and Euler-like vector fields, that are mainly used in Section~\ref{locstr}.
The {\em Euler vector field} associated to a vector bundle $E\to M$, usually denoted by $\mathcal{E}\in \mathfrak{X}(E)$, is the vector field generated by the one-parameter group $s\mapsto \kappa_{e^{-s}}$, where the maps $\kappa_{t}\colon E\to E$ are given by $\kappa_{t}e=t\cdot e$, $t\in \R$ (in case $t=0$, $\kappa_0$ is the projection to the zero section).


Given a submanifold $N$ of $M$, we denote the normal bundle of $N$ by
\[\nu(M,N)=TM|_{N}/TN,
\]
 when the context is clear we denote $\nu(M,N)$ by $\nu_N$ and the projection map by $p\colon \nu(M,N)\to N$.

Any map of pairs $\varphi\colon (M',N')\to (M,N)$ has an associated map $\nu(\varphi)\colon \nu(M',N')\to \nu(M,N).$
The following lemma is a known fact.
\begin{Lemma}\label{fiberwise_isom}
Let $\varphi\colon M'\to M$ be a smooth map. If \smash{$N\xhookrightarrow{\iota}M$} is a submanifold transverse to $\varphi$ and $N'=\varphi^{-1}(N)$, then $\nu(\varphi)$ is a fiberwise isomorphism.
\end{Lemma}
Consider a vector field $X\in \mathfrak{X}(M)$ tangent to $N$, that is $X\colon (M,N)\to (TM,TN)$. Using the fact that $T\nu_N=\nu_{TN}$ (see \cite[Appendix A]{bursztyn2019splitting}), we have that $\nu(X)$ is a vector field on~$\nu(M,N)$.
\begin{Definition}\label{defTN-EL}
Let $N$ be a submanifold of $M$. A {\em tubular neighborhood embedding} is an embedding
$\psi\colon \nu_N\to M$ that takes $N$ to $N$ and $\nu(\psi)={\rm Id}$, where $\nu(\psi)$ is induced by $\psi\colon (\nu_N,N)\to (M,N)$ (here we are using the identification $\nu(\nu_N,N)=\nu_N$).
A vector field $X\in\mathfrak{X}(M)$ is called {\em Euler-like} (along $N$) if it is complete, $X|_{N}=0$ and $\nu(X)$ is the Euler vector field of $\nu(M,N)$.
\end{Definition}
Given a submanifold $N$ of $M$, there is a one-to-one correspondence between Euler-like vector and tubular neighborhood embedding, see, for example, \cite{bursztyn2019splitting}.
\begin{Lemma}[{\cite[Lemma 3.9]{bursztyn2019splitting}}]\label{existence_esp_section}
Let $(A, \rho)$ be RAVB over $M$, and $N\subseteq M$ a submanifolds transversal to $\rho$.
Then there exists a section $\epsilon\in \Gamma(A)$ with $\epsilon|_N = 0$, such that $X = \rho(\epsilon)$ is Euler-like along~$N$.
\end{Lemma}

Given a vector bundle \smash{$E\xrightarrow{p} M$}, we know that
\[
TE\xrightarrow{Tp} TM \qquad \text{and}\qquad TE\times_E TE\xrightarrow{Tp\times Tp} TM \times _M TM
\]
 are vector bundles. We give the structure of complex vector bundle to the last one and we call it {\em the complexified tangent bundle} \smash{$T_{\C}E\xrightarrow{T_{\C}p} T_{\C}M$}, with the fiber sum and fiber multiplication by scalars given by
\begin{align*}
&(V+{\rm i} W)+_{T_{\C}p} \bigl(V'+{\rm i} W'\bigr)=\bigl(V+_{Tp} V'\bigr)+{\rm i}\bigl(W+_{Tp} W'\bigr),\\
&(\lambda_1+{\rm i}\lambda_2)\cdot_{T_{\C}p} (V+{\rm i}W)=(\lambda_1\cdot_{Tp} V-\lambda_2\cdot_{Tp} W) + {\rm i}(\lambda_2\cdot_{Tp} V + \lambda_1\cdot_{Tp} W),
\end{align*}
where $V$, $V'$, $W$ and $W'\in T_{\C}E$, with $p(V)=p(V')$, $p(W)=p(W')$ and $\lambda_1,\lambda_2\in \R$, and~$+_{Tp}$ and~$\cdot_{Tp}$ denotes the operations of fiber sum and multiplication by scalar in the tangent vector bundle \smash{$TE\xrightarrow{Tp}TM$}. Note that the previous operations make \smash{$T_{\C}E\xrightarrow{T_{\C}p} T_{\C}M$} a complex vector~bundle,
\[
\begin{tikzcd}
\Gamma\bigl(\wedge^k T^*_{\C}M\bigr)\arrow[d, "\widetilde{d}"] \arrow[r, "\widehat{}"]
& \Gamma\bigl(\wedge^k T^*M\bigr)_{\C}\arrow[d, "d_{\C}"] \\
\Gamma\bigl(\wedge^{k+1} T^*_{\C}M\bigr)\arrow[r, "\widehat{}"]
& \Gamma\bigl(\wedge^{k+1} T^*M\bigr)_{\C}.
\end{tikzcd}
\]
